% Pregunta 503 (2014-septiembre-OEP2014, n.º 35): con préstamo y endeudamiento ilimitados a R_f, la frontera eficiente es la CML, recta creciente desde R_f tangente a la frontera de Markowitz en M. Autor: Víctor Gutiérrez Marcos
\documentclass[tikz,border=6pt]{standalone}
\input{tikz-tcee}
\begin{document}
\begin{tikzpicture}[tcee]
  \ejes{$\sigma$}{$E(R)$}{6.2}{5.4}
  % Frontera de Markowitz: sigma = sqrt(4 + 5 (E - 2.2)^2); rama inferior (ineficiente) fina
  \draw[thick,tceeazul,dashed,smooth,variable=\y,domain=1.0:2.2,samples=40]
    plot ({sqrt(4+5*(\y-2.2)^2)},\y);
  % Rama superior (eficiente sin activo seguro)
  \draw[curva2,smooth,variable=\y,domain=2.2:4.2,samples=60]
    plot ({sqrt(4+5*(\y-2.2)^2)},\y) node[etiqueta,right,align=left] {frontera de\\Markowitz};
  % CML: E = 1.2 + 0.6708 sigma, tangente en M = (2.683, 3.0)
  \draw[curva1] (0,1.2) -- (5.4,4.822) node[etiqueta,above] {CML};
  \node[punto] at (0,1.2) {};
  \node[left] at (0,1.2) {$R_f$};
  \node[punto] at (2.683,3.0) {};
  \node[etiqueta,below right] at (2.683,3.0) {$M$};
  \proyeccion{(2.683,3.0)}{$\sigma_M$}{$E(R_M)$}
  % Tramos
  \node[etiqueta,tceemorado,rotate=33.85,above] at (1.3,2.07) {prestar};
  \node[etiqueta,tceemorado,rotate=33.85,above] at (4.05,3.92) {endeudarse};
  % Pendiente
  \node[etiqueta,tceemorado,anchor=west] at (0.25,4.95) {pendiente $=\dfrac{E(R_M)-R_f}{\sigma_M}>0$};
\end{tikzpicture}
\end{document}
